%--------------------------------------
% Caso de Uso: Editar Empleado
%--------------------------------------

\begin{UseCase}{CU24}{Editar Empleado}{
	Permitir que el Administrador modifique la información de un empleado previamente registrado en el sistema del Hotel para Perros.
}
	\UCitem{Versión}{\color{Gray}1.1}
	\UCitem{Autor}{\color{Gray}Equipo de Desarrollo}
	\UCitem{Supervisa}{\color{Gray}Coordinador del Sistema}
	\UCitem{Actor}{Administrador}
	\UCitem{Propósito}{Actualizar los datos personales y laborales de un empleado.}
	\UCitem{Entradas}{Nombre, email, teléfono, puesto y salario.}
	\UCitem{Origen}{Teclado}
	\UCitem{Salidas}{Datos del empleado actualizados y mensaje de confirmación.}
	\UCitem{Destino}{Pantalla de lista de empleados}
	\UCitem{Precondiciones}{
		El Administrador debe haber iniciado sesión y el empleado debe existir en el sistema.
	}
	\UCitem{Postcondiciones}{
		La información del empleado queda actualizada en la base de datos.
	}
	\UCitem{Errores}{
		Empleado no encontrado, datos inválidos, permisos insuficientes.
	}
	\UCitem{Tipo}{Caso de uso primario}
	\UCitem{Observaciones}{
		La edición solo puede ser realizada por usuarios con rol de administrador.
	}
\end{UseCase}

%--------------------------------------
\begin{UCtrayectoria}
	\UCpaso[\UCactor] Selecciona la opción \IUbutton{Editar} correspondiente a un empleado en la Pantalla de Lista de Empleados.
	\UCpaso El sistema muestra la Pantalla de Edición de Empleado con los datos actuales del empleado.
	\UCpaso[\UCactor] Modifica uno o más campos del formulario.
	\UCpaso[\UCactor] Presiona el botón \IUbutton{Guardar}.
	\UCpaso El sistema valida que los datos ingresados sean correctos.
	\UCpaso El sistema verifica que el empleado exista en el sistema.
	\UCpaso El sistema actualiza la información del empleado en la base de datos.
	\UCpaso El sistema muestra el mensaje ``Empleado actualizado correctamente''.
	\UCpaso El sistema regresa a la Pantalla de Lista de Empleados.
\end{UCtrayectoria}

%--------------------------------------
\begin{UCtrayectoriaA}{A}{Empleado no encontrado}
	\UCpaso El sistema muestra el mensaje ``El empleado seleccionado no existe''.
	\UCpaso Termina el caso de uso.
\end{UCtrayectoriaA}

%--------------------------------------
\begin{UCtrayectoriaA}{B}{Datos inválidos}
	\UCpaso El sistema muestra mensajes indicando los campos que contienen errores.
	\UCpaso Continúa en el paso 3 de la trayectoria principal.
\end{UCtrayectoriaA}

%--------------------------------------
\begin{UCtrayectoriaA}{C}{Permisos insuficientes}
	\UCpaso El sistema muestra el mensaje ``No cuenta con permisos para editar empleados''.
	\UCpaso Termina el caso de uso.
\end{UCtrayectoriaA}

%--------------------------------------
\subsection{Puntos de extensión}
\UCExtenssionPoint{
	El Administrador desea consultar el historial del empleado.
}{
	Del paso 2 al paso 4.
}{
	Consultar Empleado.
}

%--------------------------------------
% Fin del Caso de Uso

